%% ma: L11-B16-0065
%% lop: 11
%% chuong: 5
%% bai: 16
%% dang: 11.16.8
%% loai: TLN
%% muc: TH
%% dap_an: "30"
%% nguon: "tai-lieu-hoc-tap-mon-toan-lop-11-theo-chuong-trinh-giao-duc-pho-thong-2018.pdf (Lớp Toán Thầy Duy, Phạm Lê Duy, THPT Chu Văn An), Chương 3 Bài 2, Phần 3, Câu 3 (tr.378)"
%% kien_thuc: "Lớp 11 Bài 16 (giới hạn của hàm số)"
%% trang_thai: chinh-thuc
%% ghi_chu: "Đã duyệt 03/11/2026."
\begin{ex}
Một cái hồ chứa $600$ l nước ngọt. Người ta bơm nước biển có nồng độ muối $30$ g/l vào hồ với tốc độ $15$ lít/phút. Nồng độ muối của nước trong hồ sau $t$ phút kể từ khi bắt đầu bơm là $C(t)=\dfrac{30\cdot15t}{600+15t}=\dfrac{30t}{40+t}$ (g/l). Khi đó nồng độ muối trong hồ sẽ bằng bao nhiêu gam/lít khi $t$ dần về dương vô cùng?
\shortans{30}
\loigiai{\sloppy $\lim\limits_{t\to+\infty}\dfrac{30t}{40+t}=\lim\dfrac{30}{\frac{40}t+1}=30$ (g/l).}
\end{ex}
